%
% cli-session.tex
%
% Z Specification: the CLI named-session lifecycle.
% A standalone state machine (live, dormant, ended) with a host pointer,
% an alias held by at most one non-ended record, and a mailbox that
% catches every send whose address no longer names a holder.
%
% Type-check: fuzz -t cli-session.tex
% Animate:    probcli cli-session.tex -animate 20
%

\documentclass[a4paper,10pt,fleqn]{article}
\usepackage[margin=1in]{geometry}
\usepackage{fuzz}
\usepackage[colorlinks=true,linkcolor=blue,citecolor=blue,urlcolor=blue]{hyperref}

\begin{document}

\title{The CLI Named Session: A Z Specification}
\author{Lifecycle, Addressing, and Delivery}
\date{November 2026}
\hypersetup{
  pdftitle={The CLI Named Session: A Z Specification},
  pdfauthor={Punt Labs},
  pdfsubject={Z Specification},
  pdfcreator={fuzz/probcli}
}
\maketitle

\tableofcontents
\newpage

% ============================================================
\section{Introduction}
% ============================================================

A session started from the command line, rather than by an engine
host, has a name that a person types: a \emph{tty alias}.  The record
behind the name passes through three states.  It is \emph{live} while
in use, \emph{dormant} after a retention lapse, and \emph{ended} once
closed.  A per-repository \emph{pointer} remembers which record the
host last used, and may point at a record that has since ended.

This specification is deliberately separate from
\texttt{session-model.tex}.  That file is the structural model of
organisations, repositories, machines, clones, and engine sessions,
with its own routing layer and its own $Invariant$ schema; the CLI
session is a second lifecycle whose state overlaps none of those
components, and bolting it on would make a large file larger than
ProB can usefully explore.

The two files are related by one correspondence.
\texttt{session-model.tex} declares
$ttyalias : SESSION \pinj ALIAS$, so that no alias names two live
entries.  Here the same property is stated over the CLI records and
this file's own $ALIAS$ set: the $tty$ field is injective over the
\emph{non-ended} records (invariant~1).  An ended record releases its
alias, so a recycled name never collides with the record that used to
hold it, and never inherits that record's mail, because a new record
takes a fresh identifier.

The eight properties owed, each stated so that ProB decides it over the
reachable state space:

\begin{enumerate}
  \item \textbf{Alias injectivity} over the non-ended records.
  \item \textbf{Id freshness}: an identifier, once minted, is never
    minted again.
  \item \textbf{Frozen identity}: $tty$, $user$, $repo$, and $created$
    of a record never change after it is made.
  \item \textbf{Pointer validity}: every pointer entry names an
    existing record of the pointer's own repository.
  \item \textbf{Ended is drained}: an ended record holds no mail.
  \item \textbf{Message conservation}: every message sent and not yet
    read is held exactly once, in an inbox or a user mailbox.
  \item \textbf{No wrong-session delivery}: under the rule, neither
    latch fires; the defect operations are reachable only under the
    defect mode, as fidelity controls.
  \item \textbf{No auto-revive}: an ended session never becomes live
    again; a name returns only through a new record.
\end{enumerate}

% ============================================================
\section{Basic Types}
% ============================================================

\begin{zed}
  [SID, ALIAS, USER, REPO, MSG]
\end{zed}

$SID$ is the identifier of a session record, minted once.  $ALIAS$ is
the tty alias a person types.  $MSG$ is a message; its content is not
modelled, only where it is held.

\begin{zed}
  TIME == 0 \upto 2
\end{zed}

$TIME$ is a bounded natural, so that ProB can enumerate it.

\begin{zed}
  STATE ::= live | dormant | ended
\end{zed}

\begin{zed}
  ZBOOL ::= ztrue | zfalse
\end{zed}

\begin{zed}
  ROUTEMODE ::= rmRule | rmDefect
\end{zed}

$ROUTEMODE$ is a scoping parameter in the manner of
\texttt{session-model.tex}.  Under $rmRule$ delivery follows the rule:
a send reaches the holder of the address at send time, or the addressed
user's mailbox when nobody holds it.  Under $rmDefect$ the defect
operations are enabled, so that a model check can show the invariants
\emph{do} fail when the rule is broken.

% ============================================================
\section{The Record and the World}
% ============================================================

\begin{schema}{CliRecord}
  tty : ALIAS \\
  user : USER \\
  repo : REPO \\
  created : TIME
\end{schema}

A record is made once and never altered.

\begin{schema}{CliWorld}
  records : SID \pfun CliRecord \\
  state : SID \pfun STATE \\
  usedids : \power SID \\
  endedids : \power SID \\
  current : REPO \pfun SID \\
  inbox : SID \pfun \power MSG \\
  usermail : USER \pfun \power MSG \\
  sentm : \power MSG \\
  readm : \power MSG \\
  routeMode : ROUTEMODE \\
  lostF : ZBOOL \\
  wrongF : ZBOOL
  \where
  \dom state = \dom records \\
  \dom inbox = \dom records \\
  \dom usermail = USER \\
  %%
  %% (2) Id freshness: every id ever minted is remembered
  %%
  \dom records \subseteq usedids \\
  endedids \subseteq usedids \\
  %%
  %% (1) Alias injectivity over the non-ended records
  %%
  \forall a, b : \dom state |
    state~a \neq ended \land state~b \neq ended
      \land (records~a).tty = (records~b).tty @ a = b \\
  %%
  %% (4) Pointer validity
  %%
  \forall r : \dom current @
    current~r \in \dom records
      \land (records~(current~r)).repo = r \\
  %%
  %% (5) Ended is drained
  %%
  \forall s : \dom state | state~s = ended @ inbox~s = \emptyset \\
  %%
  %% (6) Message conservation: what was sent and not read is held,
  %% and held exactly once
  %%
  readm \subseteq sentm \\
  \forall m : MSG @
    m \in sentm \setminus readm \iff
      ((\exists s : \dom inbox @ m \in inbox~s)
        \lor (\exists u : \dom usermail @ m \in usermail~u)) \\
  \forall s, t : \dom inbox | s \neq t @ inbox~s \cap inbox~t = \emptyset \\
  \forall u, v : \dom usermail | u \neq v @
    usermail~u \cap usermail~v = \emptyset \\
  \forall s : \dom inbox; u : \dom usermail @
    inbox~s \cap usermail~u = \emptyset \\
  %%
  %% (8) No auto-revive: an id that has ended is ended wherever it
  %% still has a record
  %%
  \forall s : endedids \cap \dom state @ state~s = ended \\
  %%
  %% (7) No wrong-session delivery: under the rule, neither latch fires
  %%
  routeMode = rmRule \implies lostF = zfalse \\
  routeMode = rmRule \implies wrongF = zfalse
\end{schema}

$records$ holds every record not yet expired, and $state$ gives each its
state.  $usedids$ is grow-only and is the memory that makes an
identifier fresh for ever; $endedids$ is the grow-only set of
identifiers that have ended or been expired, kept so that invariant~8
can be a state predicate rather than a transition one.  $current$ is
the host pointer.  Mail is held as sets of messages rather than counts,
because a count cannot say that a message is held \emph{once}: $sentm$
and $readm$ are ghost histories of messages sent and read, and
invariant~6 equates the held mail with their difference.  $lostF$ and
$wrongF$ are history latches, as in \texttt{session-model.tex}.  No
operation but a defect operation sets either, and no operation clears
either.

ProB selects $CliWorld$ through $Init$ and requires the reserved name
$Invariant$ to install the predicates as the machine-checkable
invariant.

\begin{schema}{Invariant}
  CliWorld
\end{schema}

% ============================================================
\section{Initial State}
% ============================================================

\begin{schema}{Init}
  CliWorld'
  \where
  records' = \emptyset \\
  state' = \emptyset \\
  usedids' = \emptyset \\
  endedids' = \emptyset \\
  current' = \emptyset \\
  inbox' = \emptyset \\
  usermail' = \{ u : USER @ u \mapsto \emptyset \} \\
  sentm' = \emptyset \\
  readm' = \emptyset \\
  lostF' = zfalse \\
  wrongF' = zfalse
\end{schema}

$routeMode'$ is left free, so ProB explores $rmRule$ and $rmDefect$ as
separate branches; no operation reassigns it.

% ============================================================
\section{Frames}
% ============================================================

Each operation names what it changes and leaves the rest to one of the
frame schemas below.  $FrozenIdentity$ is the transition form of
invariant~3, a two-state predicate that every operation includes: a
record present before and after is the same record.

\begin{schema}{FrozenIdentity}
  \Delta CliWorld
  \where
  \forall s : \dom records \cap \dom records' @ records'~s = records~s
\end{schema}

\begin{schema}{KeepFlags}
  \Delta CliWorld
  \where
  routeMode' = routeMode \\
  lostF' = lostF \\
  wrongF' = wrongF
\end{schema}

\begin{schema}{KeepHistory}
  \Delta CliWorld
  \where
  sentm' = sentm \\
  readm' = readm
\end{schema}

\begin{schema}{KeepIds}
  \Delta CliWorld
  \where
  usedids' = usedids \\
  endedids' = endedids
\end{schema}

\begin{schema}{KeepMail}
  \Delta CliWorld
  \where
  inbox' = inbox \\
  usermail' = usermail
\end{schema}

\begin{schema}{KeepRecords}
  \Delta CliWorld
  \where
  records' = records \\
  state' = state
\end{schema}

\begin{schema}{KeepPointer}
  \Delta CliWorld
  \where
  current' = current
\end{schema}

\begin{schema}{Holds}
  CliWorld \\
  u? : USER \\
  a? : ALIAS
  \where
  \exists h : \dom state @
    state~h \neq ended
      \land (records~h).tty = a?
      \land (records~h).user = u?
\end{schema}

$Holds$ is the test, not an operation: some non-ended record of user
$u?$ holds alias $a?$.  By invariant~1 there is then exactly one.

% ============================================================
\section{Lifecycle Operations}
% ============================================================

\subsection{New}

A new record takes an identifier outside $usedids$ and an alias that no
non-ended record holds, starts $live$, has an empty inbox, and becomes
the pointer for its repository.

\begin{schema}{New}
  \Delta CliWorld \\
  FrozenIdentity \\
  KeepFlags \\
  KeepHistory \\
  s? : SID \\
  a? : ALIAS \\
  u? : USER \\
  r? : REPO \\
  t? : TIME
  \where
  s? \notin usedids \\
  \lnot (\exists h : \dom state @
    state~h \neq ended \land (records~h).tty = a?) \\
  (\exists rec : CliRecord |
    rec.tty = a? \land rec.user = u? \land rec.repo = r?
      \land rec.created = t? @
    records' = records \oplus \{ s? \mapsto rec \}) \\
  state' = state \oplus \{ s? \mapsto live \} \\
  usedids' = usedids \cup \{ s? \} \\
  endedids' = endedids \\
  current' = current \oplus \{ r? \mapsto s? \} \\
  inbox' = inbox \oplus \{ s? \mapsto \emptyset \} \\
  usermail' = usermail
\end{schema}

\subsection{Attach}

Attaching a $live$ record leaves it live; attaching a $dormant$ record
makes it live.  An $ended$ or absent record is not enabled: the
precondition is $state~s? \in \{live, dormant\}$, so nothing revives.

\begin{schema}{Attach}
  \Delta CliWorld \\
  FrozenIdentity \\
  KeepFlags \\
  KeepHistory \\
  KeepIds \\
  KeepMail \\
  KeepPointer \\
  s? : SID
  \where
  s? \in \dom state \\
  state~s? \in \{ live, dormant \} \\
  records' = records \\
  state' = state \oplus \{ s? \mapsto live \}
\end{schema}

\subsection{Lapse}

\begin{schema}{Lapse}
  \Delta CliWorld \\
  FrozenIdentity \\
  KeepFlags \\
  KeepHistory \\
  KeepIds \\
  KeepMail \\
  KeepPointer \\
  s? : SID
  \where
  s? \in \dom state \\
  state~s? = live \\
  records' = records \\
  state' = state \oplus \{ s? \mapsto dormant \}
\end{schema}

\subsection{Use}

Use is enabled for a non-ended record only, and makes it the pointer of
its own repository.  It writes no ended target, which is what lets
$Select$ treat an ended target as stale.

\begin{schema}{Use}
  \Delta CliWorld \\
  FrozenIdentity \\
  KeepFlags \\
  KeepHistory \\
  KeepIds \\
  KeepMail \\
  KeepRecords \\
  s? : SID
  \where
  s? \in \dom state \\
  state~s? \neq ended \\
  current' = current \oplus \{ (records~s?).repo \mapsto s? \}
\end{schema}

\subsection{End}

Ending drains the inbox into the owner's mailbox, releases the alias
(by leaving the non-ended set, which is all injectivity quantifies
over), and sets $ended$.  The pointer is left alone.

\begin{schema}{End}
  \Delta CliWorld \\
  FrozenIdentity \\
  KeepFlags \\
  KeepHistory \\
  KeepPointer \\
  s? : SID
  \where
  s? \in \dom state \\
  state~s? \neq ended \\
  records' = records \\
  state' = state \oplus \{ s? \mapsto ended \} \\
  usermail' = usermail \oplus
    \{ (records~s?).user \mapsto
         usermail~((records~s?).user) \cup inbox~s? \} \\
  inbox' = inbox \oplus \{ s? \mapsto \emptyset \} \\
  usedids' = usedids \\
  endedids' = endedids \cup \{ s? \}
\end{schema}

\subsection{Expire}

A retention lapse removes a $dormant$ or $ended$ record.  Its mail is
drained first, and any pointer to it is cleared, so that invariant~4
holds.  Its identifier stays in $usedids$.

\begin{schema}{Expire}
  \Delta CliWorld \\
  FrozenIdentity \\
  KeepFlags \\
  KeepHistory \\
  s? : SID
  \where
  s? \in \dom state \\
  state~s? \in \{ dormant, ended \} \\
  records' = \{ s? \} \ndres records \\
  state' = \{ s? \} \ndres state \\
  usermail' = usermail \oplus
    \{ (records~s?).user \mapsto
         usermail~((records~s?).user) \cup inbox~s? \} \\
  inbox' = \{ s? \} \ndres inbox \\
  current' = current \nrres \{ s? \} \\
  usedids' = usedids \\
  endedids' = endedids \cup \{ s? \}
\end{schema}

\subsection{Select}

Selection resolves a repository's pointer.  A pointer whose target is
not ended is kept as it stands.  A stale pointer, meaning an ended
target or none, falls through to $New$.

\begin{schema}{SelectHit}
  \Xi CliWorld \\
  r? : REPO
  \where
  r? \in \dom current \\
  state~(current~r?) \neq ended
\end{schema}

\begin{schema}{SelectStale}
  New
  \where
  r? \notin \dom current \lor state~(current~r?) = ended
\end{schema}

% ============================================================
\section{Delivery}
% ============================================================

A message is addressed to a user and an alias.  Three operations
implement the rule; two more are the defect.

\subsection{SendToHolder}

\begin{schema}{SendToHolder}
  \Delta CliWorld \\
  FrozenIdentity \\
  KeepFlags \\
  KeepIds \\
  KeepRecords \\
  KeepPointer \\
  h? : SID \\
  u? : USER \\
  a? : ALIAS \\
  m? : MSG
  \where
  h? \in \dom state \\
  state~h? \neq ended \\
  (records~h?).tty = a? \\
  (records~h?).user = u? \\
  m? \notin sentm \\
  inbox' = inbox \oplus \{ h? \mapsto inbox~h? \cup \{ m? \} \} \\
  usermail' = usermail \\
  sentm' = sentm \cup \{ m? \} \\
  readm' = readm
\end{schema}

\subsection{SendStaleToMailbox}

When nobody holds the address, the message goes to the addressed user's
mailbox, never to another user and never to a session that does not hold
the alias.

\begin{schema}{SendStaleToMailbox}
  \Delta CliWorld \\
  FrozenIdentity \\
  KeepFlags \\
  KeepIds \\
  KeepRecords \\
  KeepPointer \\
  u? : USER \\
  a? : ALIAS \\
  m? : MSG
  \where
  \lnot (\exists h : \dom state @
    state~h \neq ended
      \land (records~h).tty = a?
      \land (records~h).user = u?) \\
  m? \notin sentm \\
  usermail' = usermail \oplus \{ u? \mapsto usermail~u? \cup \{ m? \} \} \\
  inbox' = inbox \\
  sentm' = sentm \cup \{ m? \} \\
  readm' = readm
\end{schema}

\subsection{ReadInbox and ReadMail}

\begin{schema}{ReadInbox}
  \Delta CliWorld \\
  FrozenIdentity \\
  KeepFlags \\
  KeepIds \\
  KeepRecords \\
  KeepPointer \\
  s? : SID \\
  m? : MSG
  \where
  s? \in \dom inbox \\
  m? \in inbox~s? \\
  inbox' = inbox \oplus \{ s? \mapsto inbox~s? \setminus \{ m? \} \} \\
  usermail' = usermail \\
  sentm' = sentm \\
  readm' = readm \cup \{ m? \}
\end{schema}

\begin{schema}{ReadMail}
  \Delta CliWorld \\
  FrozenIdentity \\
  KeepFlags \\
  KeepIds \\
  KeepRecords \\
  KeepPointer \\
  u? : USER \\
  m? : MSG
  \where
  m? \in usermail~u? \\
  usermail' = usermail \oplus
    \{ u? \mapsto usermail~u? \setminus \{ m? \} \} \\
  inbox' = inbox \\
  sentm' = sentm \\
  readm' = readm \cup \{ m? \}
\end{schema}

\subsection{SendMisroute (defect)}

The defect: a send whose address names no holder is routed to a session
$w?$ that does not hold it.  The operation only latches $wrongF$ and
moves nothing, since what matters is that it fires.  It is enabled
under $rmDefect$ alone.  With $rmDefect$ chosen, invariant~7 fails at
the first firing, which is the evidence that the check is not vacuous.

\begin{schema}{SendMisroute}
  \Delta CliWorld \\
  FrozenIdentity \\
  KeepHistory \\
  KeepIds \\
  KeepMail \\
  KeepRecords \\
  KeepPointer \\
  u? : USER \\
  a? : ALIAS \\
  w? : SID
  \where
  routeMode = rmDefect \\
  w? \in \dom state \\
  \lnot (\exists h : \dom state @
    state~h \neq ended
      \land (records~h).tty = a?
      \land (records~h).user = u?) \\
  wrongF' = ztrue \\
  lostF' = lostF \\
  routeMode' = routeMode
\end{schema}

\subsection{SendLost (defect)}

The second control: a send to a stale address that is dropped rather
than mailed.  It latches $lostF$ and is likewise enabled only under
$rmDefect$.

\begin{schema}{SendLost}
  \Delta CliWorld \\
  FrozenIdentity \\
  KeepHistory \\
  KeepIds \\
  KeepMail \\
  KeepRecords \\
  KeepPointer \\
  u? : USER \\
  a? : ALIAS
  \where
  routeMode = rmDefect \\
  \lnot (\exists h : \dom state @
    state~h \neq ended
      \land (records~h).tty = a?
      \land (records~h).user = u?) \\
  lostF' = ztrue \\
  wrongF' = wrongF \\
  routeMode' = routeMode
\end{schema}

% ============================================================
\section{Verification}
% ============================================================

\begin{verbatim}
fuzz -t docs/cli-session.tex
probcli docs/cli-session.tex -model_check -strict \
  -card SID 2 -card ALIAS 2 -card USER 1 -card REPO 1 -card MSG 1 \
  -p DEFAULT_SETSIZE 1 -p MAXINT 2 -p MAX_OPERATIONS 300000 -coverage
\end{verbatim}

The primary run is exhaustive at two identifiers, two aliases, one user,
one repository, and one message: 5514 states analysed, no invariant
violation, no deadlock.  Two identifiers suffice to exercise a recycled alias (end one,
then $New$ the other on the same alias).  Because the invariant is
stated as an implication on $routeMode$, it holds in both branches; the
question the controls answer is whether the latches can fire at all:

\begin{verbatim}
probcli ... -goal "wrongF=ztrue & routeMode=rmDefect"
probcli ... -goal "lostF=ztrue & routeMode=rmDefect"
probcli ... -goal "(wrongF=ztrue or lostF=ztrue) & routeMode=rmRule"
\end{verbatim}

The first two goals are found, so the check is not vacuous: the latches
are reachable when the defect is selected.  The third reports no such
state after visiting all 5514 states: under the rule neither latch can
fire.

\end{document}
